\documentclass[11pt]{article}

\usepackage[letterpaper,margin=1in]{geometry}
\usepackage[T1]{fontenc}
\usepackage[utf8]{inputenc}
\usepackage{lmodern}
\usepackage{microtype}
\usepackage{amsmath,amssymb,mathtools,bm}
\usepackage{enumitem}
\usepackage{xcolor}
\usepackage{booktabs}
\usepackage{array}
\usepackage{hyperref}
\usepackage[most]{tcolorbox}
\usepackage{fancyhdr}
\usepackage{titlesec}
\usepackage{parskip}
\usepackage{needspace}

\definecolor{qftblue}{HTML}{245A8D}
\definecolor{qftgreen}{HTML}{2F6B4F}
\definecolor{qftred}{HTML}{A23B3B}
\definecolor{qftgray}{HTML}{F4F6F8}

\hypersetup{
  colorlinks=true,
  linkcolor=qftblue,
  urlcolor=qftblue,
  citecolor=qftgreen,
  pdftitle={PHYS 662 Lecture 10: Projective Representations and Noether Charges},
  pdfauthor={Nima Lashkari}
}

\pagestyle{fancy}
\fancyhf{}
\lhead{PHYS 662: Quantum Field Theory I}
\rhead{Lecture 10}
\cfoot{\thepage}
\setlength{\headheight}{14pt}

\titleformat{\section}{\Large\bfseries\color{qftblue}}{\thesection.}{0.6em}{}
\titleformat{\subsection}{\large\bfseries\color{qftgreen}}{\thesubsection}{0.6em}{}
\titleformat{\subsubsection}{\normalsize\bfseries\color{qftred}}{\thesubsubsection}{0.6em}{}

\setlist[itemize]{leftmargin=2em,itemsep=0.25em,topsep=0.4em}
\setlist[enumerate]{leftmargin=2em,itemsep=0.25em,topsep=0.4em}

\newcommand{\dd}{\mathrm{d}}
\newcommand{\R}{\mathbb{R}}
\newcommand{\Lag}{\mathcal{L}}
\newcommand{\cM}{\mathcal{M}}
\newcommand{\SO}{\mathrm{SO}}
\newcommand{\OO}{\mathrm{O}}

\newtcolorbox{keyidea}[1][]{
  enhanced, breakable, colback=qftgray, colframe=qftblue,
  boxrule=0.8pt, arc=1.5mm, title={#1}, fonttitle=\bfseries,
  coltitle=white
}
\newtcolorbox{remarkbox}[1][]{
  enhanced, breakable, colback=white, colframe=qftgreen,
  boxrule=0.7pt, arc=1.5mm, title={#1}, fonttitle=\bfseries,
  coltitle=white
}
\newtcolorbox{warningbox}[1][]{
  enhanced, breakable, colback=white, colframe=qftred,
  boxrule=0.7pt, arc=1.5mm, title={#1}, fonttitle=\bfseries,
  coltitle=white
}


\newcommand{\Th}{\Theta}
\newcommand{\Eul}{\mathcal E}
\newcommand{\Lie}{\mathscr L}

\newcommand{\C}{\mathbb C}
\newcommand{\Z}{\mathbb Z}
\newcommand{\SU}{\mathrm{SU}}
\newcommand{\UU}{\mathrm{U}}
\newcommand{\GL}{\mathrm{GL}}
\newcommand{\SL}{\mathrm{SL}}
\newcommand{\Sp}{\mathrm{Sp}}
\newcommand{\tr}{\operatorname{tr}}
\newcommand{\Ad}{\operatorname{Ad}}
\newcommand{\ad}{\operatorname{ad}}
\newcommand{\lieg}{\mathfrak g}
\newcommand{\lieh}{\mathfrak h}
\newcommand{\su}{\mathfrak{su}}
\newcommand{\so}{\mathfrak{so}}
\newcommand{\slie}{\mathfrak{sl}}


\newcommand{\cH}{\mathcal H}
\newcommand{\Spin}{\operatorname{Spin}}
\newcommand{\id}{\mathrm{id}}
\begin{document}
\pagenumbering{gobble}
\begin{titlepage}
\centering
\vspace*{1.2in}
{\scshape\Large PHYS 662 --- Quantum Field Theory I\par}
\vspace{0.35in}
{\Huge\bfseries\color{qftblue} Lecture 10\par}
\vspace{0.15in}
{\LARGE\bfseries Projective Representations\par}
\vspace{0.06in}
{\LARGE\bfseries and Noether Charges\par}
\vspace{0.16in}
{\large\color{qftgreen} Central Extensions and the Algebra of Symmetries\par}
\vspace{0.6in}
{\large Nima Lashkari\par}
{\large Fall 2026\par}
\vfill
\begin{minipage}{0.86\textwidth}
\small
Quantum symmetries can compose up to a phase. We derive the constraints
on that phase and its infinitesimal central term, contrasting rotational
spin with the Galilean mass extension. We then return to fields: a
complex scalar multiplet supplies explicit Noether currents, canonical
generators, and a complete calculation of their charge algebra.

\medskip
Course page: \url{https://nimalashkari.com/quantum-field-theory-2026-part-i/}
\end{minipage}
\vspace{0.5in}
\end{titlepage}
\pagenumbering{roman}
\begingroup\small
\tableofcontents
\endgroup
\newpage
\pagenumbering{arabic}

\section*{Roadmap}
\addcontentsline{toc}{section}{Roadmap}
Lecture 9 explained why rotations can act on quantum rays even when
their vector implementers compose only up to phase. We now ask which
phases can be removed and what remains in the generator algebra. Then
we derive that algebra directly from fields. The goals are:
\begin{enumerate}
\item derive the associativity and rephasing rules for projective phases;
\item distinguish a global covering-group effect from a nonremovable
central Lie algebra term;
\item verify conservation, generator action, and charge brackets for a
complex scalar multiplet.
\end{enumerate}
The ideals and semisimple algebras of Lecture 8 enter when we ask which
central terms can be removed. The field representations of Lecture 9
enter when we construct currents.

\begin{keyidea}[Conventions and two kinds of generators]
We set \(\hbar=1\), use mostly-minus Minkowski signature, and write
\(U(\theta)=e^{-i\theta^aT_a}\). The finite matrices \(t_a\) acting on
field components obey \([t_a,t_b]=if_{ab}{}^ct_c\). The corresponding
charges \(Q_a\) act on the full field phase space or quantum Hilbert
space; they are different objects from those component matrices.

Our canonical convention is \(\{q,p\}=1\) and
\(\delta F=\epsilon^a\{F,Q_a\}\). At the quantum level the same
transformation is \(U^\dagger F U\), with
\(U=e^{-i\epsilon^aQ_a}\). Keeping the order of each bracket fixed will
make the signs in the current and charge algebra agree.
\end{keyidea}

\section{Projective multiplication and phase freedom}
\label{sec:projective}
\subsection{Associativity constrains the phase}
Restrict for now to unitary implementers. Two successive transformations
agree with their group product on rays, so on vectors they may obey
\begin{equation}
 \boxed{U(g)U(h)=e^{i\omega(g,h)}U(gh).}
 \label{eq:projective-law}
\end{equation}
This is a projective representation. The phase is independent of the
state vector: two linear implementers of the same ray transformation
differ by a scalar, not by an arbitrary state-dependent phase.

Apply the projective law twice to each association:
\begin{align*}
 (U(g)U(h))U(k)&=e^{i[\omega(g,h)+\omega(gh,k)]}U(ghk),\\
 U(g)(U(h)U(k))&=e^{i[\omega(h,k)+\omega(g,hk)]}U(ghk).
\end{align*}
The resulting operators must agree. Thus associativity gives
\begin{equation}
 \omega(g,h)+\omega(gh,k)
 =\omega(h,k)+\omega(g,hk)\pmod{2\pi}.
 \label{eq:group-cocycle}
\end{equation}
This is the \emph{cocycle condition}. Choosing \(U(e)=I\) normalizes
\(e^{i\omega(e,g)}=e^{i\omega(g,e)}=1\). Near the identity we may choose
a smooth real branch with \(\omega(e,g)=\omega(g,e)=0\).
A globally smooth choice of implementers or real phase function need not
exist; all Taylor expansions below are local.

If antiunitary elements are included, moving a scalar past \(U(g)\)
complex conjugates it when \(U(g)\) is antiunitary. The right side of
Eq.~\eqref{eq:group-cocycle} becomes
\(s(g)\omega(h,k)+\omega(g,hk)\), where \(s(g)=+1\) for unitary and
\(-1\) for antiunitary implementers. We use the ordinary, untwisted
condition only on the unitary part of the symmetry group.

\subsection{Rephasing and central extensions of groups}
Changing the phase convention of each implementer,
\(U'(g)=e^{i\alpha(g)}U(g)\), changes the multiplier to
\begin{equation}
 \omega'(g,h)=\omega(g,h)+\alpha(g)+\alpha(h)-\alpha(gh)
 \pmod{2\pi}.
 \label{eq:coboundary}
\end{equation}
Such changes do not alter the action on rays. A multiplier is removable
if a single-valued choice of phases makes the representation law exact.
Local removability does not guarantee a globally consistent choice.

\begin{remarkbox}[A multiplier that is only a phase convention]
Start with an ordinary representation of the additive real line,
\(R(t)=e^{-itP}\), and choose \(U(t)=e^{i\kappa t^2}R(t)\), with
\(\kappa\) real. Then
\[
 U(s)U(t)=e^{i\kappa(s^2+t^2-(s+t)^2)}U(s+t)
         =e^{-2i\kappa st}U(s+t).
\]
The nonzero multiplier disappears upon rephasing by
\(e^{-i\kappa t^2}\). A nonzero phase in one formula is therefore
not enough to establish a nontrivial extension.
\end{remarkbox}

One can absorb a multiplier into an enlarged group. Choose one
implementer \(U(g)\) for each g in a neighborhood (a local
\emph{section}). Use pairs \((g,z)\), \(z\in\UU(1)\), with multiplication
\begin{equation}
 (g,z)(h,w)=\bigl(gh,zw e^{i\omega(g,h)}\bigr).
 \label{eq:extension-product}
\end{equation}
Equation~\eqref{eq:group-cocycle} is exactly associativity of this product.
The phase subgroup is central, and
\(\widetilde U(g,z)=zU(g)\) obeys an ordinary representation law.
The phase coordinate can be chosen locally; around the whole group,
different choices may have to be patched together.

\begin{keyidea}[Two questions that must be kept separate]
\textbf{Global question:} Can the phases be chosen consistently around
noncontractible loops in the group? Covering groups detect this issue.

\textbf{Infinitesimal question:} Do the generator commutators contain a
central term that no constant shift of generators removes?
Lie algebra central extensions detect this issue.

A projective representation can have an ordinary Lie algebra. The
half-integer-spin representations of \(\SO(3)\) in Lecture 9 gave an
explicit example.
\end{keyidea}

\Needspace{10\baselineskip}
\section{Central terms in the generator algebra}
\label{sec:central}
We now look close to the identity. Lecture 8 obtained an ordinary Lie
bracket by comparing two orders of composition. Repeating that argument
with the phase retained shows precisely where a central term can enter.

\subsection{Differentiating the projective law}
Choose local exponential coordinates near the identity. The group law is
\[
 F^c(\theta,\varphi)=\theta^c+\varphi^c+
 \tfrac12 f_{ab}{}^c\theta^a\varphi^b+O(3),
\]
and a normalized smooth multiplier begins as
\begin{equation}
 \omega(\theta,\varphi)=\chi_{ab}\theta^a\varphi^b+O(3).
 \label{eq:local-multiplier}
\end{equation}
Pure terms in only one parameter vanish by normalization. Write
\(U(\theta)=I-i\theta^aT_a+\tfrac12\theta^a\theta^bC_{ab}+O(3)\),
where \(C_{ab}=C_{ba}\). Comparing the mixed terms in
\(U(\theta)U(\varphi)=e^{i\omega}U(F)\) gives
\[
 -T_aT_b=C_{ab}-\tfrac i2 f_{ab}{}^cT_c+i\chi_{ab}I.
\]
Subtract the equation with a and b exchanged:
\begin{equation}
 \boxed{[T_a,T_b]=if_{ab}{}^cT_c+iz_{ab}I},
 \qquad z_{ab}=\chi_{ba}-\chi_{ab}=-z_{ba}.
 \label{eq:central-algebra}
\end{equation}
The antisymmetric part of the multiplier supplies the extra commutator
term; its symmetric part cancels. For Hermitian generators,
\(z_{ab}\) is real.

The identity commutes with every operator. Abstractly, one adjoins a
central generator Z, writes
\([T_a,T_b]=if_{ab}{}^cT_c+iz_{ab}Z\), and represents Z by I in the
chosen sector. More generally there can be several central generators,
whose scalar values label irreducible sectors.

\subsection{Jacobi identities and removable terms}
The Jacobi identity imposes both the usual constraint on f and
\begin{equation}
 f_{bc}{}^d z_{ad}+f_{ca}{}^d z_{bd}+f_{ab}{}^d z_{cd}=0.
 \label{eq:central-jacobi}
\end{equation}
Thus one cannot append arbitrary constants to a Lie algebra. This is the
infinitesimal cocycle condition.

A shift \(T'_a=T_a+b_aI\) does not change the induced ray transformation,
but changes the central coefficients. The left side of the commutator
is unchanged; on the right substitute \(T_c=T'_c-b_cI\):
\[
 [T'_a,T'_b]=if_{ab}{}^cT'_c+i(z_{ab}-f_{ab}{}^cb_c)I.
\]
Hence
\begin{equation}
 z'_{ab}=z_{ab}-f_{ab}{}^c b_c.
 \label{eq:central-shift}
\end{equation}
A term of the form \(z_{ab}=f_{ab}{}^cb_c\) can therefore be removed.
The space of allowed central terms modulo such shifts is denoted
\(H^2(\lieg,\R)\). For this lecture, the cocycle equation and shift
formula are the operational meaning of that notation.

For example, an apparent central term in the rotation algebra is always
of the form \(z_{ij}=\epsilon_{ijk}b_k\), because every antisymmetric
three-by-three matrix has this form. Shifting \(J_i\) to \(J_i+b_iI\)
removes it. This direct calculation already shows why the spin-one-half
representation is not explained by a nontrivial central extension of
\(\so(3)\).

\Needspace{12\baselineskip}
\subsection{What semisimplicity does imply}
The relevance of semisimplicity to central extensions is a further
structural result:

\begin{keyidea}[Whitehead's second lemma, in the form we need]
For a finite-dimensional semisimple Lie algebra over \(\R\) or \(\C\),
every scalar central Lie algebra extension is removable. Equivalently,
\(H^2(\lieg,\R)=0\) for a real semisimple algebra with trivial scalar
coefficients. A trivial direct-sum central extension can always be added;
the theorem rules out a nontrivial extension class.
\end{keyidea}
This is a theorem about the algebra, not a claim that every global group
with that algebra has only ordinary representations. Passing to a simply
connected cover removes the remaining topological obstruction in the
usual continuous unitary representation setting. On SO(3) itself,
half-integer spins still give projective representations.

Nor does the theorem follow just from tracelessness of finite-dimensional
semisimple generators. Quantum charges may act on infinite-dimensional
spaces, and tracing their commutators is generally unjustified.
The cohomological statement has the correct scope; a proof is given in
the further reading.

\subsection{The Poincar\'e example and its dimension restriction}
The Poincar\'e algebra is not semisimple: translations form a nonzero
abelian ideal. Nevertheless in spacetime dimension \(D>2\), scalar
central terms in the ordinary Poincar\'e algebra are removable. In
3+1 dimensions the connected group's universal cover is
\(\R^{1,3}\rtimes\SL(2,\C)\). Its ordinary unitary representations
include the representations that are projective on the proper
orthochronous Poincar\'e group, so the absence of a nontrivial Lie algebra
central term does not exclude half-integer spin.

The dimensional qualification matters. In 1+1 dimensions, let H, P, and
K generate time translations, spatial translations, and boosts. The
extension
\begin{equation}
 [K,H]=iP,\qquad [K,P]=iH,\qquad [H,P]=iZ,
 \qquad [Z,\cdot]=0
 \label{eq:poincare-two-dimensional}
\end{equation}
satisfies Jacobi. Its last bracket cannot be removed by constant shifts
of H and P, because their unextended bracket is zero. We will use the
3+1-dimensional result for relativistic QFT; the lower-dimensional
example shows why its dimension restriction is necessary.

\section{A nontrivial extension: mass in Galilean symmetry}
\label{sec:galilei}
The Galilei group provides a physical central extension that survives
infinitesimally. Consider a free nonrelativistic particle with mass m,
\(\{x_i,p_j\}=\delta_{ij}\), and
\begin{equation}
 H=\frac{\bm p^2}{2m},\qquad P_i=p_i,\qquad K_i=mx_i-tp_i.
 \label{eq:galilean-generators}
\end{equation}
With our choice of boost parameter,
\(\delta x_i=\{x_i,v_jK_j\}=-tv_i\) and
\(\delta p_i=-mv_i\). Reversing the boost parameter gives the opposite
coordinate convention used in some texts. Direct calculation yields
\begin{equation}
 \{K_i,P_j\}=m\delta_{ij},\qquad
 \{K_i,H\}=P_i,\qquad \{K_i,K_j\}=0.
 \label{eq:galilean-poisson}
\end{equation}
The explicitly time-dependent boost charge is nevertheless conserved:
\(\dd K_i/\dd t=\partial_tK_i+\{K_i,H\}=-p_i+p_i=0\).
In particular, conservation need not mean that the charge separately
Poisson-commutes with H when it has explicit time dependence.

Quantization gives \([K_i,P_j]=im\delta_{ij}I\). Abstractly m is the
value of a central mass operator M. For several particles it is the total
mass. A shift of K or P by a constant does not change their bracket, so
this extension cannot be removed.

There is also a direct group-level test. Set
\(B(\bm v)=e^{-i\bm v\cdot\bm K}\) and
\(T(\bm a)=e^{-i\bm a\cdot\bm P}\). Since the commutator of their
exponents is central, the Baker--Campbell--Hausdorff formula gives exactly
\begin{equation}
 B(\bm v)T(\bm a)
 =e^{-im\bm v\cdot\bm a}T(\bm a)B(\bm v).
 \label{eq:bargmann-phase}
\end{equation}
Boosts and spatial translations commute as transformations of Galilean
spacetime, but their quantum implementers differ by this phase. The
extended group is often called the Bargmann group. Unlike the SU(2)
double cover of SO(3), this example adds a nontrivial continuous central
direction to the algebra itself. Notice also that the central term was
already present in the classical Poisson algebra: it need not be a
quantum anomaly.

\section{Noether currents for a complex scalar multiplet}
\label{sec:currents}
We now connect the representation matrices to quantities built from
fields. There are three steps: derive a conserved current, check that
its charge generates the field transformation, and compute the brackets
of two charges. Conservation alone does not perform the last two checks.

\subsection{The symmetry and its equations of motion}
Return to relativistic fields in \(d+1\) dimensions. Let \(\Phi\) be a
column of N complex scalar fields, with Lagrangian density
\begin{equation}
 \Lag=\partial_\mu\Phi^\dagger\partial^\mu\Phi-V(s),
 \qquad s=\Phi^\dagger\Phi.
 \label{eq:multiplet-lagrangian}
\end{equation}
For a differentiable real potential V, this is invariant under a constant
internal transformation
\begin{equation}
 \Phi'=U\Phi,\qquad \Phi'^\dagger=\Phi^\dagger U^\dagger,
 \qquad U\in\UU(N).
 \label{eq:multiplet-action}
\end{equation}
The fundamental and dual/conjugate actions and their invariant pairing
were developed in Lecture 9. Here they ensure that both the
kinetic term and s are invariant.
Space-dependent U would produce extra derivative terms and would require
a different theory with a gauge connection.
Written as \(2N\) real scalars, this particular Lagrangian also has an
O(2N) symmetry; we study its chosen U(N) subgroup and the associated
\(N^2\) currents.

Choose Hermitian matrices \(t_a\), \(a=0,\ldots,N^2-1\), normalized by
\begin{equation}
 \tr(t_at_b)=\tfrac12\delta_{ab},\qquad
 t_0=\frac{I_N}{\sqrt{2N}},\qquad [t_a,t_b]=if_{ab}{}^ct_c.
 \label{eq:uN-basis}
\end{equation}
There is one central u(1) generator and \(N^2-1\) traceless su(N)
generators. Treat \(\Phi\) and \(\Phi^\dagger\) as independent variables
when varying the action. The resulting equations are
\begin{equation}
 \Box\Phi+V'(s)\Phi=0,\qquad
 \Box\Phi^\dagger+V'(s)\Phi^\dagger=0.
 \label{eq:multiplet-eom}
\end{equation}

\subsection{Deriving and checking the currents}
For \(U=e^{-i\epsilon^at_a}\),
\(\delta\Phi=-i\epsilon^at_a\Phi\) and
\(\delta\Phi^\dagger=i\epsilon^a\Phi^\dagger t_a\).
There is no spacetime displacement or total-derivative term in this
strict internal symmetry. The Noether formula gives
\begin{align}
 j_a^\mu
 &=\partial^\mu\Phi^\dagger(-it_a\Phi)
   +(i\Phi^\dagger t_a)\partial^\mu\Phi\notag\\
 &=\boxed{i\bigl(\Phi^\dagger t_a\partial^\mu\Phi
       -\partial^\mu\Phi^\dagger t_a\Phi\bigr)}.
 \label{eq:noether-current}
\end{align}
These are \(N^2\) real classical currents, one for each independent
symmetry parameter.
On shell,
\begin{align}
 \partial_\mu j_a^\mu
 &=i\bigl(\Phi^\dagger t_a\Box\Phi
     -\Box\Phi^\dagger t_a\Phi\bigr)\notag\\
 &=i\bigl(-V'(s)\Phi^\dagger t_a\Phi
          +V'(s)\Phi^\dagger t_a\Phi\bigr)=0.
 \label{eq:current-divergence}
\end{align}
The first-derivative terms cancel before using the equations of motion.
The remaining terms vanish by the equations of motion.

The charge is the integral of the time component,
\begin{equation}
 Q_a=\int\dd^dx\,j_a^0
 =i\int\dd^dx\,
 \bigl(\Phi^\dagger t_a\dot\Phi
       -\dot\Phi^\dagger t_a\Phi\bigr).
 \label{eq:noether-charge}
\end{equation}
Its time derivative is minus the boundary flux
\(\int_{\partial\Sigma}\dd S_i\,j_a^i\); assume this flux vanishes.
Reversing the parameter in the field transformation reverses both the
current and charge together.

\subsection{How the currents transform}
Define \(D(U)\) by \(Ut_bU^\dagger=t_aD(U)^a{}_b\).
Trace invariance makes D orthogonal in the basis
\eqref{eq:uN-basis}. Inserting \(\Phi'=U\Phi\) in the current gives
\begin{equation}
 j_a'^\mu=D(U^{-1})^b{}_a j_b^\mu
          =D(U)^a{}_b j_b^\mu.
 \label{eq:current-adjoint}
\end{equation}
Thus the currents carry the coadjoint representation, identified with
the adjoint by the invariant trace pairing. The matrix
\(\mathcal J^\mu=2j_a^\mu t_a\) transforms as
\(\mathcal J'^\mu=U\mathcal J^\mu U^\dagger\).
As explained in Lecture 9, the U(N) adjoint is reducible: its u(1)
component is a singlet and its traceless part is the SU(N) adjoint.
The central u(1) acts trivially by conjugation, even though its Noether
charge generates a nontrivial phase rotation of the fundamental field.

\section{The Poisson algebra and quantization of charges}
\label{sec:charges}
\subsection{Canonical variables and the generator test}
At a fixed time, the canonical momenta are
\begin{equation}
 \Pi_A=\frac{\partial\Lag}{\partial\dot\Phi^A}
       =\dot\Phi_A^\dagger,
 \qquad
 \overline\Pi^{\,A}=\frac{\partial\Lag}{\partial\dot\Phi_A^\dagger}
       =\dot\Phi^A.
 \label{eq:momenta}
\end{equation}
Capital indices label field components; lower-case a,b label generators.
In complex canonical coordinates the nonzero elementary brackets are
\begin{equation}
 \{\Phi^A(\bm x),\Pi_B(\bm y)\}
 =\delta^A{}_B\delta^{(d)}(\bm x-\bm y),\qquad
 \{\Phi_A^\dagger(\bm x),\overline\Pi^{\,B}(\bm y)\}
 =\delta_A{}^B\delta^{(d)}(\bm x-\bm y).
 \label{eq:canonical-pb}
\end{equation}
The reverse-ordered brackets have minus signs. All cross brackets vanish.
These complex coordinates are a convenient rewriting of the usual real
canonical coordinates; they do not introduce extra degrees of freedom.

In matrix notation, with \(\Pi\) a row and \(\overline\Pi\) a column,
\begin{equation}
 Q_a=i\int\dd^dx\,
       \bigl(\Phi^\dagger t_a\overline\Pi-\Pi t_a\Phi\bigr).
 \label{eq:charge-canonical}
\end{equation}
In the bracket with \(\Phi\), only the term containing \(\Pi\)
contributes. Its delta function collapses the spatial integral:
\begin{align}
 \{\Phi^A(\bm y),Q_a\}
 &=-i\int\dd^dx\,
 \delta^A{}_B\delta^{(d)}(\bm y-\bm x)
 (t_a)^B{}_C\Phi^C(\bm x)\notag\\
 &=-i(t_a\Phi)^A(\bm y),\label{eq:field-generator}\\
 \{\Phi_A^\dagger(\bm y),Q_a\}
 &=i(\Phi^\dagger t_a)_A(\bm y).
 \label{eq:conjugate-generator}
\end{align}
Thus \(\delta F=\epsilon^a\{F,Q_a\}\) reproduces precisely the
transformation used to derive the current. If the charge is placed first
in the bracket, its action changes sign.

\begin{remarkbox}[One complex field is a rotation of two real fields]
For a concrete check take N=1 and use the phase parameterization
\(\delta\phi=-i\epsilon\phi\), with generator 1 rather than the
trace-normalized \(t_0\). Write
\(\phi=(q_1+iq_2)/\sqrt2\). Its real components transform as
\[
 \delta q_1=\epsilon q_2,\qquad \delta q_2=-\epsilon q_1.
\]
With \(p_i=\dot q_i\), the current and charge become
\[
 j^\mu=q_2\partial^\mu q_1-q_1\partial^\mu q_2,
 \qquad Q=\int\dd^dx\,(q_2p_1-q_1p_2).
\]
The real canonical brackets give \(\{q_1,Q\}=q_2\) and
\(\{q_2,Q\}=-q_1\), exactly as required. The charge is an angular
momentum in the internal \((q_1,q_2)\) plane; the sign corresponds
to the orientation chosen by \(e^{-i\epsilon}\).
\end{remarkbox}

\subsection{Computing the charge brackets}
A short bilinear identity keeps the matrix multiplication visible. For
constant matrices A and B define
\[
 F[A]=\int\dd^dx\,\Pi A\Phi,\qquad
 G[A]=\int\dd^dx\,\Phi^\dagger A\overline\Pi.
\]
For the first bracket, the functional derivatives are
\(\delta F[A]/\delta\Phi^I=(\Pi A)_I\) and
\(\delta F[A]/\delta\Pi_I=(A\Phi)^I\). Thus the canonical rule gives
\begin{align*}
 \{F[A],F[B]\}
 &=\int\dd^dx\,\bigl[(\Pi A)_I(B\Phi)^I
                      -(A\Phi)^I(\Pi B)_I\bigr]\\
 &=\int\dd^dx\,\Pi(AB-BA)\Phi=F([A,B]).
\end{align*}
Repeating this for the conjugate pair yields the three useful identities
\begin{equation}
 \{F[A],F[B]\}=F([A,B]),\qquad
 \{G[A],G[B]\}=-G([A,B]),\qquad
 \{F[A],G[B]\}=0.
 \label{eq:bilinear-brackets}
\end{equation}
The opposite signs come from the opposite order of coordinate and
momentum indices in the two bilinears. The mixed bracket vanishes
because the two canonical pairs have zero cross brackets.
Since \(Q_a=-iF[t_a]+iG[t_a]\),
\begin{align}
 \{Q_a,Q_b\}
 &=-F([t_a,t_b])+G([t_a,t_b])\notag\\
 &=f_{ab}{}^c\bigl(-iF[t_c]+iG[t_c]\bigr)
 =\boxed{f_{ab}{}^cQ_c}.
 \label{eq:charge-algebra}
\end{align}
The abstract Lie algebra is therefore realized by functions on the full
classical field phase space. It is not merely an algebra satisfied by the
finite matrices \(t_a\).

As a check, the u(1) charge \(Q_0\) commutes with every \(Q_a\), because
\(f_{0a}{}^b=0\). Yet \(\{\Phi,Q_0\}=-i\Phi/\sqrt{2N}\) is nonzero.
Being central among symmetry charges does not mean acting trivially on
all fields.

\subsection{Where a central term could have entered}
Suppose charges generate the expected symmetry algebra on all canonical
variables. The Jacobi identity then implies that
\[
 C_{ab}=\{Q_a,Q_b\}-f_{ab}{}^cQ_c
\]
Poisson-commutes with all those variables. On a connected nondegenerate
canonical phase space it must be a constant. Generator actions by
themselves therefore determine the charge algebra only up to a possible
central term. The explicit calculation \eqref{eq:charge-algebra} proves
that it vanishes for this multiplet with the chosen normalization.
The Galilean example showed that this constant need not vanish in a
different theory.

Constant shifts change C as in Eq.~\eqref{eq:central-shift}, so the same
test for removability applies classically. Boundary contributions can
also affect the charge algebra in systems with different boundary
conditions.

\subsection{Quantization and the factor of i}
Canonical quantization replaces the elementary Poisson brackets by
commutators according to
\begin{equation}
 [\widehat\Phi^A(\bm x),\widehat\Pi_B(\bm y)]
 =i\delta^A{}_B\delta^{(d)}(\bm x-\bm y),
 \qquad \{\ ,\ \}_{\rm PB}\ \longrightarrow\ \frac1i[\ ,\ ].
 \label{eq:quantization}
\end{equation}
Restoring units, the denominator is \(i\hbar\). This correspondence
applies to the elementary variables and present charges, but need not
hold for arbitrary classical functions. With symmetry-preserving
regularization and operator ordering, the
scalar multiplet's global charges satisfy
\begin{equation}
 [\widehat Q_a,\widehat Q_b]=if_{ab}{}^c\widehat Q_c,
 \qquad [\widehat Q_a,\widehat\Phi]=-t_a\widehat\Phi.
 \label{eq:quantum-charge-algebra}
\end{equation}
Consequently
\[
 e^{i\epsilon^a\widehat Q_a}\widehat\Phi
 e^{-i\epsilon^a\widehat Q_a}
 =e^{-i\epsilon^at_a}\widehat\Phi,
\]
which matches the classical component transformation. Local field
products require a regulator; unbounded charges need a common invariant
domain. The stated assumptions give the unextended scalar charge
algebra. Other theories may have central terms or quantum anomalies
(failure of a classical symmetry after quantization). These are separate
questions requiring a quantum calculation.

For a conserved charge with explicit time dependence, the quantum
conservation equation is
\(\partial_t\widehat Q+i[\widehat H,\widehat Q]=0\).
For a time-independent internal charge it reduces to
\([\widehat H,\widehat Q]=0\), paralleling the classical discussion.



\clearpage
\section*{Questions and checks}
\addcontentsline{toc}{section}{Questions and checks}
\begin{enumerate}
\item Derive the cocycle condition by associating three unitary
implementers in two ways. Derive the rephasing formula as well.
\item Show that the product \eqref{eq:extension-product} is associative
exactly when the cocycle condition holds. Identify its central subgroup.
\item Derive the local central term \eqref{eq:central-algebra}, its
Jacobi constraint, and its change under a constant generator shift.
\item Remove an arbitrary antisymmetric central term in the rotation
algebra explicitly. Explain why this does not remove half-integer spin
as a projective representation of SO(3).
\item Compute \(\{K_i,P_j\}\), \(\{K_i,H\}\), and \(\dd K_i/\dd t\)
for the free particle. Derive the boost-translation phase and explain why
constant shifts cannot remove the mass term.
\item Starting from \eqref{eq:multiplet-lagrangian}, derive the currents
and verify their divergence using the field equations. Count the U(N)
and SU(N) currents separately.
\item Prove both bilinear bracket identities in
Eq.~\eqref{eq:bilinear-brackets} in component notation. Use them to
recover \(\{Q_a,Q_b\}=f_{ab}{}^cQ_c\).
\item Explain how a central u(1) charge can act nontrivially on a
fundamental field while the same central subgroup acts trivially in the
adjoint representation. Track the factor of i after quantization.
\end{enumerate}

\clearpage
\section*{Summary}
\addcontentsline{toc}{section}{Summary}
\begin{keyidea}[The completed connection]
The group law determines the infinitesimal algebra. Quantum rays allow
phases in the group law. Global topology and central Lie algebra terms
are two distinct sources of projective behavior. Noether charges realize
the symmetry as canonical transformations, and their Poisson brackets
reproduce the Lie algebra, possibly with a central term. Quantization
then introduces the factor of i in the commutator algebra.
\end{keyidea}
\begin{itemize}
\item Associativity constrains a projective multiplier. Rephasing changes
it by a coboundary; the removable-multiplier example shows why a nonzero
phase alone does not establish a nontrivial extension.
\item A central Lie algebra term satisfies Jacobi and is defined modulo
constant shifts of generators. For a finite-dimensional semisimple
algebra every scalar central extension is removable.
\item The covering-group origin of spin in Lecture 9 leaves the rotation
Lie algebra unchanged. Galilean mass instead supplies a nonremovable
central term, already present in the classical Poisson brackets.
\item For a complex scalar multiplet,
\(j_a^\mu=i(\Phi^\dagger t_a\partial^\mu\Phi
-\partial^\mu\Phi^\dagger t_a\Phi)\). The field equations imply
current conservation, and vanishing boundary flux gives conserved charges.
\item The canonical brackets verify \(\{\Phi,Q_a\}=-it_a\Phi\)
and \(\{Q_a,Q_b\}=f_{ab}{}^cQ_c\). The generator test and the
charge-algebra calculation are separate steps.
\item With the stated quantum assumptions,
\([\widehat Q_a,\widehat Q_b]=if_{ab}{}^c\widehat Q_c\).
Central terms and anomalies require attention to the actual theory,
regularization, and boundary conditions.
\end{itemize}

\section*{Further reading}
\addcontentsline{toc}{section}{Further reading}
These notes reorganize the supplied symmetry and representation slide
sets into Lectures 8--10. For mathematical background and further proofs:
\begin{itemize}
\item J. C. Fiallo R., \href{https://www.math.ubc.ca/~reichst/Lie-Algebra-Cohomology.pdf}{\emph{Lie Algebra Cohomology}},
Corollary 5.13, for Whitehead's second lemma.
\item Y. Nakayama, \emph{Central Extension of Scaling Poincar\'e Algebra},
\href{https://arxiv.org/abs/2311.10254}{arXiv:2311.10254},
for the dimension dependence of Poincar\'e central extensions.
\end{itemize}
\end{document}
